\DocumentMetadata{
  pdfversion=2.0,
  pdfstandard=ua-2,
  lang=en-US,
  tagging=on,
  tagging-setup={math/setup=mathml-SE,tabsorder=structure}
}
\documentclass[11pt,letterpaper]{article}
\usepackage[margin=0.68in,includefoot]{geometry}
\usepackage{newcomputermodern}
\usepackage{amsmath,amscd,mathtools}
\usepackage{tikz,adjustbox}
\providecommand{\square}{\mdlgwhtsquare}
\usepackage[hidelinks]{hyperref}
\hypersetup{
  pdftitle={Probability PhD Exam, May 2003},
  pdfauthor={University of Florida Department of Mathematics},
  pdfsubject={Graduate mathematics examination},
  pdfdisplaydoctitle=true,
  bookmarksopen=true
}
\setcounter{secnumdepth}{0}
\setlength{\parindent}{0pt}
\setlength{\parskip}{0.28em}
\setlength{\emergencystretch}{3em}
\setlength{\leftmargini}{2.15em}
\setlength{\leftmarginii}{2.65em}
\setlength{\leftmarginiii}{3em}
\pagestyle{empty}
\raggedbottom
\allowdisplaybreaks
\NewTaggingSocketPlug{block/list/label}{examproblem}{%
  \tagstructbegin{tag=Lbl}\tagstructend%
  \tagstructbegin{tag=\LBody}%
  \tagstructbegin{tag=\ProblemHeadingTag,title-o={\ProblemHeadingText}}%
  \tagmcbegin{tag=\ProblemHeadingTag,actualtext={\ProblemHeadingText}}%
  #2%
  \tagmcend%
  \tagstructend%
}
\newenvironment{examproblems}[2]{%
  \def\ProblemHeadingTag{#2}%
  \AssignTaggingSocketPlug{block/list/label}{examproblem}%
  \begin{enumerate}%
  \setcounter{enumi}{#1}%
  \renewcommand{\labelenumi}{\textbf{\theenumi.}}%
  \setlength{\topsep}{0.35em}%
  \setlength{\partopsep}{0pt}%
  \setlength{\itemsep}{0.48em}%
  \setlength{\parsep}{0.18em}%
}{\end{enumerate}}
\newenvironment{examsubparts}{%
  \AssignTaggingSocketPlug{block/list/label}{default}%
  \begin{enumerate}%
  \setlength{\topsep}{0.18em}%
  \setlength{\partopsep}{0pt}%
  \setlength{\itemsep}{0.16em}%
  \setlength{\parsep}{0.1em}%
}{\end{enumerate}}
\newenvironment{examinstructions}{%
  \par\begingroup\itshape
}{%
  \par\endgroup\vspace{0.18em}
}
\makeatletter
\renewcommand\subsection{\@startsection{subsection}{2}{0pt}%
  {-1.25ex plus -.25ex minus -.1ex}%
  {0.55ex plus .1ex}%
  {\normalfont\large\bfseries}}
\renewcommand{\@maketitle}{%
  \newpage\null\vskip -1em%
  {\footnotesize\bfseries
    \href{https://www.ufl.edu/}{University of Florida}, \href{https://math.ufl.edu/}{Department of Mathematics}\par}%
  \vskip -0.5em%
  \tagstructbegin{tag=H1,title={Probability PhD Exam, May 2003}}%
  \tagpdfparaOff%
  \tagmcbegin{tag=H1}%
    {\LARGE\bfseries\href{https://gma.math.ufl.edu/exams/probability-phd/}{Probability PhD Exam}, May 2003\par}%
  \tagmcend%
  \tagpdfparaOn%
  \tagstructend%
  \vskip 0.25em%
  \tagmcbegin{artifact}%
    \hrule height 1pt%
  \tagmcend%
  \vskip 1em%
}
\makeatother
\title{Probability PhD Exam, May 2003}
\author{}
\date{}
\begin{document}
\maketitle
\thispagestyle{empty}
\begin{examinstructions}
Write clearly. State precisely results you use. Do not omit steps.
\end{examinstructions}
\begin{examproblems}{0}{H2}
\def\ProblemHeadingText{Problem 1.}
\item
Let \((\overline X_m)_{m=1}^{\infty}\) be an i.i.d. sequence of real valued variables on \((\Omega,\mathcal F,P)\). Let \(\mathcal F_m=\sigma(\overline X_1,\ldots,\overline X_m)\). Put \(S_m=\overline X_1+\cdots+\overline X_m\), for \(m\geq1\), and put \(S_0=0\). Which of the following sequences are martingales relative to \(\mathcal F_m\)? For those which are not martingales with respect to this filtration, what conditions have to be added to make them martingales?
\begin{examsubparts}
\item[\textnormal{a)}]
\(S_m\), \(m=0,1,2,\ldots\), \(E[|\overline X_1|]<\infty\).
\item[\textnormal{b)}]
\(\overline X_1^2+\cdots+\overline X_m^2-m\lambda\), \(m=1,2,\ldots\), \(E(\overline X_1^2)<\infty\), \(\lambda\) real.
\item[\textnormal{c)}]
\(\exp(\alpha S_m-m\lambda)\), \(m=0,1,2,\ldots\), where \(\varphi(\alpha)=E[\exp(\alpha\overline X_1)]<\infty\), \(\alpha,\lambda\) real.
\item[\textnormal{d)}]
\(\overline Y_m=S_{\min(m,T)}\), \(m=0,1,2,\ldots\), where \(P(\overline X_1=\pm1)=\frac12\) and \(T=\min\{n>0:S_n=0\}\).
\end{examsubparts}
\def\ProblemHeadingText{Problem 2.}
\item
Let \(\{\alpha_m\}_{m=1,2,\ldots}\) be a sequence of probability measures on \(\mathbb R\). Show the following are equivalent:
\begin{examsubparts}
\item[\textnormal{a)}]
There exists a probability measure~\(\alpha\) on \(\mathbb R\) such that \(\lim_{m\to\infty}\alpha_m(I)=\alpha(I)\) for all closed intervals~\(I\) on \(\mathbb R\) whose end points are continuity points of~\(\alpha\).
\item[\textnormal{b)}]
If \(\{F_m\}\), respectively~\(F\), are the distribution functions of \(\{\alpha_m\}\), respectively~\(\alpha\), then \(\lim_{m\to\infty}F_m(x)=F(x)\) at every point~\(x\) of continuity of~\(F\).
\end{examsubparts}
\def\ProblemHeadingText{Problem 3.}
\item
Let \(\{\alpha_m\}_{m=1,2,\ldots}\) be a sequence of probability measures on \(\mathbb R\), and let \(\{\varphi_m(t)\}_{m=1,2,\ldots}\) be the corresponding characteristic functions. Assume that \(\lim_{m\to\infty}\varphi_m(t)=\varphi(t)\) for all \(t\in\mathbb R\), and \(\varphi(t)\) is continuous at \(t=0\). Prove that \(\varphi(t)\) is the characteristic function of some probability measure and that the conditions of Problem 2 hold.
\def\ProblemHeadingText{Problem 4.}
\item
If \(\overline X\) is a real valued random variable, show that for \(0<\varepsilon<1\),
\[
E\left[\frac{|\overline X|}{1+|\overline X|}\right]\leq \varepsilon+P[|\overline X|>\varepsilon]
\]
 and that
\[
P(|\overline X|>\varepsilon)\leq\frac{1+\varepsilon}{\varepsilon}E\left[\frac{|\overline X|}{1+|\overline X|}\right].
\]
 Deduce that the set of real valued random variables on a given probability space can (with suitable identification) be regarded as a metric space if the distance between two random variables \(\overline X\) and \(\overline Y\) is defined to be
\[
d(\overline X,\overline Y)=E\left[\frac{|\overline X-\overline Y|}{1+|\overline X-\overline Y|}\right],
\]
 and that convergence in this metric space corresponds to convergence in probability.
\def\ProblemHeadingText{Problem 5.}
\item
This problem has three parts.
\begin{examsubparts}
\item[\textnormal{a)}]
Let \(\overline X_m\), \(m=1,2,\ldots\), be i.i.d., \(\overline X_m>0\), \(E[\overline X_m]=1\), and \(\overline Y_m=\prod_{k=1}^m\overline X_k\). Prove that \(\overline Y_m\) is a martingale relative to the filtration \(\mathcal F_m=\sigma(\overline X_1,\ldots,\overline X_m)\). Show that \(\lim_{m\to\infty}\overline Y_m=\overline Y_\infty\) exists a.s.
\item[\textnormal{b)}]
Show that \(P[\overline Y_\infty=0]=0\) or~\(1\).

\par
Hint: Recall infinite products and the zero-one law.
\item[\textnormal{c)}]
Using b), show that
\[
E[\overline Y_\infty]=1\quad\Longleftrightarrow\quad P[\overline Y_\infty>0]=1.
\]
\end{examsubparts}
\def\ProblemHeadingText{Problem 6.}
\item
Let \(\overline X_m\), \(m=1,2,\ldots\), be a supermartingale such that
\[
E[|\overline X_{m+1}-\overline X_m|\mid\mathcal F_m]<M\quad\text{a.s.}
\]
 Show that for any two optional times \(S\leq T\), \(E[T]<\infty\), we have \(E[|\overline X_T|]<\infty\) and
\[
E(\overline X_T\mid\mathcal F_S)\leq\overline X_S.
\]
\end{examproblems}
\end{document}
